% 题干转写：完成
% 解析转写：完成
% 图形核对：完成
\section{第38届复赛第6题：同轴带电圆筒与电磁角动量}

\subsection{题干}
真空中有两个同轴、无限长、均匀带电的薄壁圆筒，半径分别为 $r_1$、$r_2$。沿轴向单位长度内，两圆筒的质量均为 $m$，内、外圆筒的电荷分别为 $q$ 和 $-q$（$q>0$），两筒均可绕轴自由旋转，如图\ref{fig:cpho38q06}。真空介电常量和磁导率分别为 $\varepsilon_0$、$\mu_0$。

求：（1）距轴 $r$ 处的电场；（2）两筒以相同角速度 $\omega$ 同向旋转时的磁场；（3）质量为 $\mu$、电荷量为 $Q$ 且 $Q$ 与 $q$ 同号的点电荷从 $r=r_1$ 处静止释放，能够到达 $r=r_2$ 的 $\omega$ 条件；（4）两筒初始静止，对外筒施加外力矩，待外筒角速度达到 $\Omega$ 时，求内筒角速度、外力矩总冲量矩和两筒总机械角动量。
\begin{figure}[htbp]
  \centering
  \begin{tikzpicture}[x=.9cm,y=.9cm,>=Stealth,font=\small]
  % Original three-dimensional coaxial-cylinder geometry.
  \begin{scope}[xshift=-2.8cm]
    \draw[thick] (0,1.55) ellipse (1.65 and .5);
    \draw[thick] (-1.65,1.55)--(-1.65,-1.15);
    \draw[thick] (1.65,1.55)--(1.65,-1.15);
    \draw[thick,dashed] (-1.65,-1.15) arc (180:360:1.65 and .5);
    \draw[thick] (1.65,-1.15) arc (0:-180:1.65 and .5);
    \draw[thick] (0,1.55) ellipse (.82 and .25);
    \draw[thick,dashed] (-.82,1.55)--(-.82,-1.15);
    \draw[thick,dashed] (.82,1.55)--(.82,-1.15);
    \draw[thick,dashed] (0,-1.15) ellipse (.82 and .25);
    \draw[densely dashed] (0,1.55)--(0,-1.15);
    \draw[->] (0,1.55)--(.82,1.55) node[midway,above] {$r_1$};
    \draw[->] (0,1.92)--(1.65,1.92) node[midway,above] {$r_2$};
    \node at (.46,.15) {$q$};
    \node at (1.28,.15) {$-q$};
    \node[left] at (-.82,.15) {$m$};
    \node[left] at (-1.65,.15) {$m$};
    \node[below] at (0,-1.8) {同轴薄壁圆筒};
  \end{scope}

  % Auxiliary cross-section for the field regions and particle motion.
  \begin{scope}[xshift=3.0cm]
    \draw[thick] (0,0) circle (1.65);
    \draw[thick] (0,0) circle (.82);
    \fill (0,0) circle (1.3pt) node[below left] {轴};
    \draw[->] (0,0)--(.82,0) node[midway,below] {$r_1$};
    \draw[->] (0,.25)--(1.65,.25) node[midway,above] {$r_2$};
    \fill (.82,0) circle (2pt);
    \node[below right,xshift=4pt,yshift=-3pt] at (.82,0) {$Q,\mu$};
    \draw[->,very thick] (.9,0)--(1.48,0) node[midway,above] {$v_r$};
    \draw[->,thick] (1.02,.75) arc (36:95:1.25) node[midway,above] {$Q\mathbf v\times\mathbf B$};
    \node[align=center] at (0,-2.2) {横截面与粒子运动\\（辅助示意）};
  \end{scope}
\end{tikzpicture}

  \caption{第38届复赛第6题同轴带电圆筒的TikZ重绘。左图对应原题图6a，右图辅助表示场区和点电荷运动。}
  \label{fig:cpho38q06}
\end{figure}

\subsection{参考解答整理}
\paragraph{（1）电场分布。}
由圆柱对称性和高斯定理，电场沿径向，
\[
\boxed{
E(r)=
\begin{cases}
\dfrac{q}{2\pi\varepsilon_0r},&r_1<r<r_2,\\[4pt]
0,&\text{其他区域}.
\end{cases}}
\]

\paragraph{（2）磁场分布。}
内筒旋转形成的轴向单位长度等效电流为 $I=q\omega/(2\pi)$。内、外筒等效电流大小相同、方向相反。以题解规定的轴向为正，
\[
\boxed{
B(r)=
\begin{cases}
-\dfrac{\mu_0q\omega}{2\pi},&r_1<r<r_2,\\[4pt]
0,&\text{其他区域}.
\end{cases}}
\]

\paragraph{（3）点电荷到达外筒的条件。}
设点电荷的径向速度为 $v_r=\dot r$，绕轴机械角动量为 $\ell$。磁场力矩满足
\[
\frac{\mathrm d\ell}{\mathrm dt}=QB r\dot r,
\qquad
\ell=\frac12QB(r^2-r_1^2),
\]
其中使用了 $\ell(r_1)=0$。电场做功给出
\[
E_k=\frac{qQ}{2\pi\varepsilon_0}\ln\frac r{r_1}
=\frac12\left(\mu v_r^2+\frac{\ell^2}{\mu r^2}\right).
\]
在 $r=r_2$ 处要求径向动能非负，代入 $|B|=\mu_0q|\omega|/(2\pi)$，得到
\[
\boxed{
\omega^2\le
\frac{16\pi\mu r_2^2}
{\varepsilon_0\mu_0^2qQ(r_2^2-r_1^2)^2}
\ln\frac{r_2}{r_1}.}
\]

\paragraph{（4）外筒驱动时的角动量传递。}
记内、外筒在自身内部产生的磁场为
\[
B_1=\frac{\mu_0q\omega}{2\pi},
\qquad
B_2=-\frac{\mu_0q\Omega}{2\pi}.
\]
变化磁场在 $r=r_1$ 处产生切向感应电场。由法拉第定律和内筒转动方程，
\[
E_1=-\frac{\mu_0qr_1}{4\pi}\frac{\mathrm d}{\mathrm dt}(\omega-\Omega),
\qquad
mr_1^2\frac{\mathrm d\omega}{\mathrm dt}=qE_1r_1.
\]
初态两筒均静止，积分得
\[
\boxed{\omega=\alpha\Omega,\qquad
\alpha\equiv\frac{\mu_0q^2}{4\pi m+\mu_0q^2}.}
\]

在 $r=r_2$ 处，感应电场对外筒的力矩为
\[
M_2=\frac{\mu_0q^2}{4\pi}
\left(r_1^2\frac{\mathrm d\omega}{\mathrm dt}
-r_2^2\frac{\mathrm d\Omega}{\mathrm dt}\right).
\]
因此感应力矩的冲量矩为
\[
J_{\text{电场}}=\frac{\mu_0q^2}{4\pi}(\alpha r_1^2-r_2^2)\Omega.
\]
外筒的机械角动量为 $L_2=mr_2^2\Omega$，由 $J+J_{\text{电场}}=L_2$ 得外力矩总冲量矩
\[
\boxed{J=\left[mr_2^2-\frac{\mu_0q^2}{4\pi}
(\alpha r_1^2-r_2^2)\right]\Omega.}
\]
两筒总机械角动量为
\[
\boxed{L=m(r_1^2\omega+r_2^2\Omega)
=m\Omega(\alpha r_1^2+r_2^2).}
\]

\subsection{转写依据}
题干依据 \texttt{CPHO38-Q06\_transcription\_v1.md}；公式依据 \texttt{CPHO38-Q06\_formula\_check\_v1.csv} 与答案第14至16页；题图依据原题图6a重绘。本题第（1）至（4）问评分分别为2分、4分、12分、22分。
